\begin{abstract}
Deep learning-based software vulnerability detection faces two coupled barriers: models fail to generalise across projects because they overfit project-specific lexical patterns, and organisations cannot pool the proprietary code that would diversify training data. We propose VulMorph-Fed, a federated learning framework built on the observation that vulnerability manifestations vary lexically across projects while their structural characteristics (morphologies) are largely shared. VulMorph-Fed combines three components: Vulnerability-Critical Subgraph Abstraction (VCSA), which maps code tokens onto a small, fully specified taxonomy of abstract operation types and learns to down-weight vulnerability-irrelevant edges; Morphology-Conditioned Federated Prototype Aggregation (MCFPA), which shares CWE-conditioned prototypes instead of model parameters under a provable per-round $\varepsilon$-differential-privacy guarantee with explicit composition accounting; and Morphology-Guided Message Passing (MGMP), which injects the aggregated global prototypes into local graph neural network training. We evaluate on project-level cross-project splits of four public C/C++ corpora (Devign, BigVul, DiverseVul, PrimeVul) against baselines organised by privacy class and matched in data, budget, backbone and input representation, reporting mean and standard deviation over repeated runs with statistical tests. At matched per-round privacy budgets, VulMorph-Fed substantially outperforms DP-FedAvg --- whose parameter-level noise collapses detection to near-chance --- while retaining most of the discrimination of non-private federated training, and communicating only $\sim$11\,KB per client per round. The abstracted representation itself transfers: parameter-averaging FL on our morphology graphs surpasses its lexical counterparts on the multi-project corpora. The complete implementation and scripts regenerating every reported number are publicly released.
\end{abstract}
